\documentclass[11pt, oneside]{article}   	% use "amsart" instead of "article" for AMSLaTeX format
\usepackage{geometry}                		% See geometry.pdf to learn the layout options. There are lots.
\geometry{letterpaper}                   		% ... or a4paper or a5paper or ...
%\geometry{landscape}                		% Activate for rotated page geometry
%\usepackage[parfill]{parskip}    		% Activate to begin paragraphs with an empty line rather than an indent
\usepackage{graphicx}				% Use pdf, png, jpg, or eps§ with pdflatex; use eps in DVI mode
									% TeX will automatically convert eps --> pdf in pdflatex
\usepackage{amssymb}
\usepackage{amsmath}
%SetFonts

%SetFonts


\title{Semi-implicit acoustics for the linear dry-air density}
\author{Michael M. Bell}
%\date{}							% Activate to display a given date or no date

\begin{document}
\maketitle

\section{Motivation}

The semi-implicit formulation (companion note) treats the vertically propagating
acoustic modes implicitly so that the model timestep is governed by advection rather
than the speed of sound. In its original form the implicit acoustic operator is
constant-coefficient \emph{because} the dry-air continuity equation is carried for the
log-density $\xi = \ln(\rho_d/\rho_0)$: its linearized continuity is simply
$\partial \xi'/\partial t = -\partial w/\partial z$, with no density weighting.

Carrying the log-density, however, breaks discrete conservation of dry-air mass: the
per-step cubic B-spline smoothing penalty acts on the \emph{nonlinear} $\xi$ and,
through Jensen's inequality, systematically lowers $\int \rho_d$ (see the companion
diagnostic note). The cure is to carry the \emph{linear} dry-air density $\rho_d$
directly, so that the smoothing preserves $\int \rho_d$ exactly. Doing so makes the
linearized continuity carry the reference density $\hat{\rho}_d(z)$, and the acoustic
operator appears to become variable-coefficient. This note shows that, under the
mass-flux change of variable $\varphi \equiv \hat{\rho}_d\, w$, the operator reduces to
the \emph{same} constant-coefficient Helmholtz problem already solved for the $\xi$
form, so no new variable-coefficient solver is required.

\section{Governing equations and reference state}

We use the notation of the Scythe documentation. The dry-air continuity equation in
conservative flux form, and the vertical momentum equation with the Ooyama (1990)
pressure-gradient decomposition, are

\begin{align}
\frac{\partial \rho_d}{\partial t} &= -\nabla \cdot \left( \rho_d \mathbf{v} \right) , \\
\frac{\partial w}{\partial t} &= -\mathbf{v}\cdot\nabla w
   - \frac{1}{\rho}\frac{\partial p}{\partial z} - g\frac{\rho'}{\rho} + \dot{F}_w ,
\end{align}

where the vertical pressure gradient is

\begin{equation}
\frac{\partial p}{\partial z} = P_s \frac{\partial s}{\partial z}
   + \rho_d P_{\rho_d}\frac{\partial \xi}{\partial z}
   + \frac{dq_v}{d\mu}P_{q_v}\frac{\partial \mu}{\partial z} ,
\qquad
P_{\rho_d} = \frac{\partial p}{\partial \rho_d}\bigg\vert_{s,q_v} .
\end{equation}

We define a hydrostatic reference state, denoted by the hat operator, and solve for
perturbations to it:

\begin{equation}
\frac{d\hat{p}}{dz} = -\hat{\rho}\,g , \qquad
\hat{\xi} = \ln(\hat{\rho}_d/\rho_0) , \qquad
\rho_d' = \rho_d - \hat{\rho}_d , \qquad
\xi' = \ln(\rho_d/\hat{\rho}_d) \approx \frac{\rho_d'}{\hat{\rho}_d} .
\end{equation}

The last relation is the linearization of $\xi'$ used throughout the acoustic
treatment.

\section{The linearized acoustic subsystem}

Splitting each tendency into an explicit part $[\,\cdot\,]^{ex}$ and the linear
acoustic part that propagates sound vertically, the dry-acoustic pair couples
$\rho_d'$ and $w$:

\begin{align}
\frac{\partial \rho_d'}{\partial t} &= \left[\frac{\partial \rho_d'}{\partial t}\right]^{ex}
   - \frac{\partial}{\partial z}\left(\hat{\rho}_d\, w\right) , \label{eq:cont} \\
\frac{\partial w}{\partial t} &= \left[\frac{\partial w}{\partial t}\right]^{ex}
   - \bar{c}_\rho\,\frac{\partial \rho_d'}{\partial z} , \label{eq:wmom}
\end{align}

where the reference acoustic coefficient is

\begin{equation}
\bar{c}_\rho \equiv \frac{\bar{P}_\xi}{\hat{\rho}_d} , \qquad
\bar{P}_\xi \equiv \left\langle \frac{\rho_d P_{\rho_d}}{\rho} \right\rangle .
\label{eq:cbar}
\end{equation}

Here $\bar{P}_\xi$ is the domain-mean sound speed squared --- the single scalar the
code stores as \texttt{Pxi\_bar}, namely $\langle P_\xi / \rho \rangle$ with
$P_\xi = \rho_d P_{\rho_d}$. It is the \emph{same} reference constant used by the $\xi$
form; the $\rho_d$ form differs only in carrying $\hat{\rho}_d(z)$ explicitly in both
\eqref{eq:cont} and \eqref{eq:wmom}.

\paragraph{Reduction to the $\xi$ form.} With $\xi' \approx \rho_d'/\hat{\rho}_d$, the
momentum source in \eqref{eq:wmom} is
\begin{equation}
-\bar{c}_\rho\,\frac{\partial \rho_d'}{\partial z}
   = -\frac{\bar{P}_\xi}{\hat{\rho}_d}\frac{\partial \rho_d'}{\partial z}
   = -\bar{P}_\xi\,\frac{\partial \xi'}{\partial z} ,
\end{equation}
which is exactly the $\xi$-form implicit term $-\bar{P}_\xi\,\partial \xi'/\partial z$.
Likewise $-\partial_z(\hat{\rho}_d w)/\hat{\rho}_d \approx -\partial_z w$ when
$\hat{\rho}_d$ is slowly varying. A constant $\hat{\rho}_d$ recovers the $\xi$ form
identically.

\section{Time discretization}

We use the same splitting as the existing solver: a third-order Adams--Bashforth (AB3)
extrapolation of the explicit parts and the AI2$^\ast$ weighting
$\left(\tfrac{5}{4}, -1, \tfrac{3}{4}\right)$ of the implicit acoustic terms, with
$\Delta\tau$ the implicit time-stepping coefficient ($\Delta\tau = \tfrac{5}{4}\Delta t$
after the startup steps). Let $\rho_d'^{\,*}$ and $w^{*}$ denote the explicitly
advanced predictors, formed by adding the AB3 explicit increment, removing the
AB-extrapolated estimate of the implicit terms, and adding their implicit weight.
Equations \eqref{eq:cont}--\eqref{eq:wmom} then give the coupled implicit update

\begin{align}
w^{n+1} &= w^{*} - \Delta\tau\,\bar{c}_\rho\,\frac{\partial \rho_d'^{\,n+1}}{\partial z} ,
   \label{eq:wupd} \\
\rho_d'^{\,n+1} &= \rho_d'^{\,*}
   - \Delta\tau\,\frac{\partial}{\partial z}\left(\hat{\rho}_d\, w^{n+1}\right) .
   \label{eq:cupd}
\end{align}

\section{The mass-flux reduction}

Introduce the vertical \textbf{mass flux} $\varphi \equiv \hat{\rho}_d\, w$. Multiplying
\eqref{eq:wupd} by $\hat{\rho}_d$ and substituting \eqref{eq:cupd},

\begin{equation}
\varphi^{n+1} = \hat{\rho}_d\, w^{*}
   - \Delta\tau\,\hat{\rho}_d\,\bar{c}_\rho\,\frac{\partial}{\partial z}
     \left[ \rho_d'^{\,*} - \Delta\tau\,\frac{\partial \varphi^{n+1}}{\partial z} \right] .
\end{equation}

The decisive simplification is that the product
$\hat{\rho}_d\,\bar{c}_\rho = \hat{\rho}_d\,(\bar{P}_\xi/\hat{\rho}_d) = \bar{P}_\xi$
is \emph{constant}. The variable reference density cancels, leaving a
constant-coefficient modified-Helmholtz problem for the mass flux:

\begin{equation}
\boxed{\;
\varphi^{n+1} - \Delta\tau^2\,\bar{P}_\xi\,\frac{\partial^2 \varphi^{n+1}}{\partial z^2}
   = \hat{\rho}_d\, w^{*} - \Delta\tau\,\bar{P}_\xi\,\frac{\partial \rho_d'^{\,*}}{\partial z} \; }
\label{eq:helmholtz}
\end{equation}

The operator $\left(1 - \Delta\tau^2\,\bar{P}_\xi\,\partial_{zz}\right)$ is
\emph{identical} to the operator of the existing $w$-form solve
(\texttt{calc\_Helmholtz\_semiimplicit\_matrix}), which assembles
$c\,\partial_{zz} - 1$ with $c = \Delta\tau^2 \bar{P}_\xi$ and homogeneous Dirichlet
boundary rows. The boundary conditions are unchanged: a rigid lid and floor impose
$w = 0$, hence $\varphi = \hat{\rho}_d\, w = 0$, the same homogeneous Dirichlet
conditions. The already-factorized matrix is therefore reused without modification, on
either the Chebyshev (RZ) or cubic B-spline (RiRk) vertical basis.

\paragraph{Recovery.} Having solved \eqref{eq:helmholtz} for $\varphi^{n+1}$, the
physical fields are recovered by

\begin{align}
w^{n+1} &= \frac{\varphi^{n+1}}{\hat{\rho}_d} , \label{eq:wrec} \\
\rho_d'^{\,n+1} &= \rho_d'^{\,*} - \Delta\tau\,\frac{\partial \varphi^{n+1}}{\partial z} .
   \label{eq:crec}
\end{align}

\paragraph{Mass conservation.} Equation \eqref{eq:crec} updates the density through the
flux divergence $\partial_z(\hat{\rho}_d w)$. Its vertical integral is a boundary flux,
\begin{equation}
\frac{d}{dt}\int_{z_b}^{z_t}\!\rho_d'\,dz
   = -\int_{z_b}^{z_t}\!\frac{\partial}{\partial z}\!\left(\hat{\rho}_d\, w\right)dz
   = -\left[\hat{\rho}_d\, w\right]_{z_b}^{z_t} = 0
\end{equation}
at rigid walls, so the implicit acoustic step conserves $\int \rho_d'$ to machine
precision. The mass-conservation property established in the explicit linear-$\rho_d$
prototype thus carries over to the semi-implicit step.

\section{Operator consistency}

The AI2$^\ast$ split requires the implicit operator and the explicit acoustic
tendencies to apply the \emph{same} discrete derivative. Accordingly, the explicit
implicit-candidate tendencies are written in the flux form using a single spline
derivative of the mass flux, matching the recovery \eqref{eq:wrec}--\eqref{eq:crec}:

\begin{align}
\left[\frac{\partial \rho_d'}{\partial t}\right]^{\mathrm{imp}}
   &= -\frac{\partial}{\partial z}\left(\hat{\rho}_d\, w\right) , \\
\left[\frac{\partial w}{\partial t}\right]^{\mathrm{imp}}
   &= -\bar{c}_\rho\,\frac{\partial \rho_d'}{\partial z}
    = -\frac{\bar{P}_\xi}{\hat{\rho}_d}\,\frac{\partial \rho_d'}{\partial z} .
\end{align}

The continuity term must be evaluated as a single derivative of the product
$\hat{\rho}_d\, w$, \emph{not} as the product-rule expansion
$\hat{\rho}_d\,\partial_z w + w\,\partial_z\hat{\rho}_d$; the two differ in the discrete
spline operators, and only the single-derivative (flux) form is consistent with the
implicit solve \eqref{eq:helmholtz}--\eqref{eq:crec}. Using the product rule would
desynchronize the explicit and implicit acoustic operators and reintroduce a small mass
drift.

\section{Summary}

Carrying the linear dry-air density $\rho_d$ in place of $\xi$ restores discrete mass
conservation, and the resulting variable-coefficient acoustic operator collapses, under
the mass-flux variable $\varphi = \hat{\rho}_d w$, to the constant-coefficient Helmholtz
problem \eqref{eq:helmholtz} already solved for the $\xi$ form. The implementation reuses
the existing factorized matrix and adds only a new adjustment routine that (i) forms the
right-hand side of \eqref{eq:helmholtz} from the predictors $\rho_d'^{\,*}$ and
$\hat{\rho}_d w^{*}$, (ii) solves for $\varphi^{n+1}$, and (iii) recovers
$w^{n+1}$ and $\rho_d'^{\,n+1}$ by \eqref{eq:wrec}--\eqref{eq:crec}. Constant
$\hat{\rho}_d$ reduces the scheme exactly to the $\xi$-form solver, which provides a
machine-precision regression test.

\end{document}
